%==================================================================
% SECCIÓN 8 — Diseño y documentación de casos de prueba
% Entrega I · Construcción de Software y Pruebas · Grupo A · 2026-2
% Equipo: Joshua Mendez, Edinson Pedroza
% Cuerpo para \input dentro del documento maestro. Reutiliza la plantilla
% de plan-completo-ejecutado.tex: columnas L{} e Y, booktabs, tabularx a
% \textwidth y \footnotesize para las tablas anchas.
%==================================================================
\FloatBarrier
\section*{8. Diseño y documentación de casos de prueba}
\phantomsection\addcontentsline{toc}{section}{8. Diseño y documentación de casos de prueba}

\subsection*{8.1 Enfoque del diseño y técnicas de caja negra}
\addcontentsline{toc}{subsection}{8.1 Enfoque del diseño y técnicas de caja negra}

El diseño de casos para Demoblaze (\url{https://www.demoblaze.com}) se apoya en las
cuatro técnicas de diseño de caja negra que exige la rúbrica: \textbf{partición de
equivalencias}, \textbf{análisis de valores límite}, \textbf{tabla de decisión} y
\textbf{arreglos ortogonales}. Como el sistema es software de terceros —sin código
fuente ni acceso al servidor o a la base de datos—, todas las técnicas operan sobre
las entradas y salidas observables desde el navegador, que es justamente el ámbito de
la caja negra.

La suite completa está documentada en el anexo (hoja de cálculo con las columnas del
formato IEEE 829 y su registro de ejecución) y reúne \textbf{60 casos ya ejecutados},
cuyos resultados reales se consolidan en la sección 9. A esos 60 casos se suman
\textbf{8 casos nuevos de arreglos ortogonales} (\texttt{CP-AO-01} \dots{}
\texttt{CP-AO-08}), que completan la cuarta técnica y se desarrollan en el apartado
8.3. Con ello la suite alcanza \textbf{68 casos diseñados} en total.

\medskip
\noindent\textbf{Nota de integridad.} Los ocho casos \texttt{CP-AO} están
\emph{diseñados} en esta sección; su ejecución y resultados se reportan en la
sección 9. En este documento no se les atribuye ningún resultado de ejecución. Su
comportamiento esperado se deriva de las reglas de negocio observadas en el sistema y
está corroborado por los casos ya ejecutados \texttt{CP-COM-04}, \texttt{CP-COM-05} y
\texttt{CP-COM-06}, que ejercitan las mismas validaciones de obligatoriedad del
formulario de compra.

\subsection*{8.2 Cobertura por técnica de diseño}
\addcontentsline{toc}{subsection}{8.2 Cobertura por técnica de diseño}

La Tabla~\ref{tab:cobertura-tecnicas} resume cuántos casos aporta cada técnica. Las
cuatro técnicas de la rúbrica quedan cubiertas: las tres primeras ya ejecutadas dentro
de los 60 casos del plan, y los arreglos ortogonales mediante los 8 casos nuevos. Se
incluyen además la transición de estados (ciclo del carrito) y un bloque de técnicas
complementarias —comprobaciones de consistencia e inspección, más los casos de
integración, rendimiento, seguridad y compatibilidad— que no derivan de una técnica
formal de diseño de entradas pero forman parte de la suite ejecutada.

\begingroup
\small
\renewcommand{\arraystretch}{1.25}
\begin{table}[htb]\centering\small
\caption{Cobertura de la suite por técnica de diseño.}\label{tab:cobertura-tecnicas}
\begin{tabularx}{\textwidth}{@{}L{4.8cm} >{\centering\arraybackslash}p{1.5cm} L{3.4cm} Y@{}}
\toprule
\textbf{Técnica de diseño} & \textbf{N.\textsuperscript{o} de casos} & \textbf{Estado} & \textbf{Observación} \\
\midrule
Partición de equivalencias & 18 & Ejecutados (§9) & Campos de texto de los cuatro formularios (válido / inválido / vacío). \\
Valores límite & 2 & Ejecutados (§9) & Longitudes y bordes de mes/año del formulario de compra. \\
Tabla de decisión & 4 & Ejecutados (§9) & Combinaciones de sesión, carrito y campos de la compra. \\
\textbf{Arreglos ortogonales} & \textbf{8} & \textbf{Ejecutados (§9): 8/8 aprobados} & \textbf{Casos nuevos \texttt{CP-AO-01..08}; interacción por pares de los 6 campos de \emph{Place Order}.} \\
\midrule
Transición de estados & 11 & Ejecutados (§9) & Ciclo del carrito: vacío, con uno, con varios, tras eliminar y tras comprar. \\
Otras complementarias & 25 & Ejecutados (§9) & Consistencia, inspección, integración, rendimiento, seguridad y compatibilidad. \\
\midrule
\textbf{Total diseñado} & \textbf{68} & \multicolumn{2}{@{}l@{}}{\footnotesize 60 ejecutados (anexo, §9) + 8 arreglos ortogonales ejecutados (§9)} \\
\bottomrule
\end{tabularx}
\end{table}
\endgroup

Las cuatro técnicas de la rúbrica —partición de equivalencias, valores límite, tabla
de decisión y arreglos ortogonales— quedan así cubiertas. Las tres primeras ya
estaban ejercitadas en los 60 casos ejecutados; los arreglos ortogonales se añaden con
los ocho casos \texttt{CP-AO} que se diseñan a continuación.

\subsection*{8.3 Arreglos ortogonales (Taguchi): diseño de los casos CP-AO}
\addcontentsline{toc}{subsection}{8.3 Arreglos ortogonales (Taguchi): diseño de los casos CP-AO}

\subsubsection*{8.3.1 Por qué esta técnica y qué aporta a la suite}

La suite ya ejecutada aplica tres de las cuatro técnicas clásicas de caja negra. Falta
la cuarta, \textbf{arreglos ortogonales}, que es la que permite cubrir de forma
sistemática las \emph{combinaciones} entre varios parámetros de entrada. El arreglo
ortogonal resuelve un problema concreto: cuando un formulario tiene varios campos, el
número de combinaciones crece de forma exponencial y probarlas todas deja de ser
viable. Taguchi demostró que la mayoría de los defectos de interacción se disparan por
la combinación de \textbf{dos} factores a la vez, no por la de seis simultáneos. Un
arreglo ortogonal garantiza que \emph{todas las parejas} de niveles queden cubiertas
con una fracción mínima de corridas; esta propiedad se llama cobertura \textbf{pairwise}
(de a pares).

\subsubsection*{8.3.2 Función bajo prueba, factores y niveles}

Al pulsar \textbf{Place Order} sobre un carrito con productos, Demoblaze abre un
formulario con el total precargado y \textbf{seis campos editables}: Name, Country,
City, Credit card, Month y Year. Al pulsar \textbf{Purchase}, el sistema decide si
acepta o rechaza la compra en función de lo que tengan esos seis campos: una salida
que depende de la combinación de varios campos de entrada, que es el escenario para el
que sirve un arreglo ortogonal.

Tomamos los seis campos como \textbf{factores} y a cada uno le asignamos \textbf{dos
niveles} (rellenado-válido / vacío), que es la partición mínima y a la vez la más
informativa para una validación de obligatoriedad (regla R5). El dominio de valores
válidos ya se cubre con partición de equivalencias y valores límite en los casos
\texttt{CP-COM-05} y \texttt{CP-COM-06}; aquí no se repite ese trabajo, sino que se
estudia la \textbf{interacción entre los seis campos}.

\begingroup
\small
\renewcommand{\arraystretch}{1.2}
\begin{table}[htb]\centering\small
\caption{Factores y niveles del arreglo ortogonal.}\label{tab:factores}
\begin{tabularx}{\textwidth}{@{}>{\centering\arraybackslash}p{1.4cm} L{2.8cm} L{5.0cm} Y@{}}
\toprule
\textbf{Factor} & \textbf{Campo} & \textbf{Nivel 1 (rellenado-válido)} & \textbf{Nivel 2 (vacío)} \\
\midrule
A & Name & \texttt{Joshua Mendez} & \emph{(campo en blanco)} \\
B & Country & \texttt{Colombia} & \emph{(campo en blanco)} \\
C & City & \texttt{Bogotá} & \emph{(campo en blanco)} \\
D & Credit card & \texttt{4111\allowbreak1111\allowbreak1111\allowbreak1111} & \emph{(campo en blanco)} \\
E & Month & \texttt{12} & \emph{(campo en blanco)} \\
F & Year & \texttt{2030} & \emph{(campo en blanco)} \\
\bottomrule
\end{tabularx}
\end{table}
\endgroup

\subsubsection*{8.3.3 La explosión combinatoria: \texorpdfstring{$2^{6}=64$}{2\^{}6=64}}

Con 6 factores y 2 niveles cada uno, la combinación exhaustiva es

\[ 2^{6} = 2\times 2\times 2\times 2\times 2\times 2 = 64 \]

\noindent es decir, \textbf{64 corridas}. Para un solo formulario de una tienda de
demostración, probado a mano sobre un entorno de producción ajeno y compartido (véanse
las restricciones del plan de pruebas), 64 corridas son inviables y, sobre todo,
innecesarias: la enorme mayoría no aporta información nueva porque repite parejas de
condiciones ya vistas.

\subsubsection*{8.3.4 Arreglo ortogonal seleccionado: \texorpdfstring{L8 ($2^{7}$)}{L8 (2\^{}7)}}

Para 6 factores de 2 niveles, el arreglo de Taguchi más pequeño que los acomoda es el
\textbf{L8 ($2^{7}$)}: 8 corridas y hasta 7 columnas de 2 niveles. Usamos 6 de sus 7
columnas (una por factor) y dejamos la séptima sin asignar. El L8 reduce las 64
corridas exhaustivas a solo \textbf{8}. La Tabla~\ref{tab:l8-canonico} muestra la
matriz estándar en forma canónica y la Tabla~\ref{tab:l8-campos} la misma matriz
traducida a los valores de cada campo.

\begingroup
\small
\renewcommand{\arraystretch}{1.2}
\begin{table}[htb]\centering\small
\caption{Arreglo L8 en forma canónica. Cada fila es una corrida; el valor es el nivel (1 o 2) que toma el factor.}\label{tab:l8-canonico}
\setlength{\tabcolsep}{6pt}
\begin{tabular}{c cccccc c}
\toprule
\textbf{Corrida} & \textbf{A} & \textbf{B} & \textbf{C} & \textbf{D} & \textbf{E} & \textbf{F} & \textbf{Col.\,7} \\
 & {\scriptsize Name} & {\scriptsize Country} & {\scriptsize City} & {\scriptsize Card} & {\scriptsize Month} & {\scriptsize Year} & {\scriptsize (s/uso)} \\
\midrule
1 & 1 & 1 & 1 & 1 & 1 & 1 & 1 \\
2 & 1 & 1 & 1 & 2 & 2 & 2 & 2 \\
3 & 1 & 2 & 2 & 1 & 1 & 2 & 2 \\
4 & 1 & 2 & 2 & 2 & 2 & 1 & 1 \\
5 & 2 & 1 & 2 & 1 & 2 & 1 & 2 \\
6 & 2 & 1 & 2 & 2 & 1 & 2 & 1 \\
7 & 2 & 2 & 1 & 1 & 2 & 2 & 1 \\
8 & 2 & 2 & 1 & 2 & 1 & 1 & 2 \\
\bottomrule
\end{tabular}
\end{table}
\endgroup

\begingroup
\small
\renewcommand{\arraystretch}{1.2}
\begin{table}[htb]\centering\small
\caption{Arreglo L8 traducido a los campos. Nivel 1 = lleno, nivel 2 = vacío.}\label{tab:l8-campos}
\setlength{\tabcolsep}{5pt}
\begin{tabular}{c cccccc}
\toprule
\textbf{Corrida} & \textbf{Name} & \textbf{Country} & \textbf{City} & \textbf{Credit card} & \textbf{Month} & \textbf{Year} \\
\midrule
1 & lleno & lleno & lleno & lleno & lleno & lleno \\
2 & lleno & lleno & lleno & \textbf{vacío} & \textbf{vacío} & \textbf{vacío} \\
3 & lleno & \textbf{vacío} & \textbf{vacío} & lleno & lleno & \textbf{vacío} \\
4 & lleno & \textbf{vacío} & \textbf{vacío} & \textbf{vacío} & \textbf{vacío} & lleno \\
5 & \textbf{vacío} & lleno & \textbf{vacío} & lleno & \textbf{vacío} & lleno \\
6 & \textbf{vacío} & lleno & \textbf{vacío} & \textbf{vacío} & lleno & \textbf{vacío} \\
7 & \textbf{vacío} & \textbf{vacío} & lleno & lleno & \textbf{vacío} & \textbf{vacío} \\
8 & \textbf{vacío} & \textbf{vacío} & lleno & \textbf{vacío} & lleno & lleno \\
\bottomrule
\end{tabular}
\end{table}
\endgroup

\subsubsection*{8.3.5 Justificación de la reducción 64 \texorpdfstring{$\rightarrow$}{->} 8 y garantía pairwise}

¿Qué se pierde al pasar de 64 a 8 corridas? Nada de la cobertura de a pares, que es
donde se concentran los defectos de interacción. Un arreglo es \emph{ortogonal} cuando,
para \textbf{cualquier par de columnas}, las cuatro combinaciones posibles de sus
niveles —(1,1), (1,2), (2,1), (2,2)— aparecen \textbf{el mismo número de veces}. En el
L8 cada una aparece exactamente \textbf{dos veces} en cada par de columnas. Con 6
factores existen $C(6,2)=15$ parejas de campos, y el L8 cubre las cuatro combinaciones
de las 15 parejas en solo 8 corridas. La Tabla~\ref{tab:pairwise} lo verifica sobre la
pareja crítica Name--Credit card, que son los dos campos obligatorios según R5.

\begingroup
\small
\renewcommand{\arraystretch}{1.2}
\begin{table}[htb]\centering\small
\caption{Verificación de cobertura de la pareja crítica Name (A) -- Credit card (D).}\label{tab:pairwise}
\begin{tabularx}{\textwidth}{@{}L{5.4cm} Y >{\centering\arraybackslash}p{1.8cm}@{}}
\toprule
\textbf{Combinación (Name, Card)} & \textbf{Corridas donde aparece} & \textbf{Veces} \\
\midrule
(lleno, lleno) & 1, 3 & 2 \\
(lleno, vacío) & 2, 4 & 2 \\
(vacío, lleno) & 5, 7 & 2 \\
(vacío, vacío) & 6, 8 & 2 \\
\bottomrule
\end{tabularx}
\end{table}
\endgroup

Lo que el L8 \textbf{no} garantiza es la cobertura de interacciones de orden superior
(tres o más campos simultáneos en una combinación específica). Esa limitación es
aceptable y deliberada: la evidencia de Taguchi y de la literatura de \emph{pairwise
testing} indica que los fallos por combinación de tres o más factores son marginales
frente a los de a pares, y cubrirlos exigiría volver hacia las 64 corridas. Se cambia
un riesgo residual pequeño por una reducción del \textbf{87,5\,\%} en el número de
casos (de 64 a 8).

\textbf{¿Por qué 8 corridas y no 6?} Si el único objetivo fuera la cobertura pairwise
mínima, un \emph{covering array} lograría los 6 factores binarios en apenas 6 corridas.
Se usan 8 porque se exige la propiedad más fuerte de \textbf{arreglo ortogonal}: no
solo que cada par de niveles aparezca al menos una vez (covering array), sino que
aparezca el mismo número de veces (cobertura pairwise \emph{uniforme}, que obliga a que
el número de corridas sea múltiplo de 4). Esa uniformidad aporta balance estadístico
entre los factores —el fundamento del diseño de experimentos de Taguchi—; las 2
corridas extra son el precio de esa garantía de equilibrio.

\subsubsection*{8.3.6 Derivación del resultado esperado (R1, R3, R5)}

El resultado esperado de cada corrida no se inventa: se deduce de las reglas de negocio
observadas en el sistema.

\begin{itemize}
  \item \textbf{R1 — No se puede completar la compra con el carrito vacío.} Todas las
        corridas comparten la precondición de carrito con al menos un producto, de modo
        que R1 queda satisfecha y no es la variable bajo estudio. Por la misma razón, la
        \textbf{sesión iniciada} se fija como precondición constante y \textbf{no entra
        como factor}: el objetivo es la interacción entre los seis campos, no el efecto
        de la autenticación (que ya prueba por separado \texttt{CP-COM-03}).
  \item \textbf{R5 — Name y Credit card son obligatorios.} Si Name o Credit card quedan
        vacíos, el sistema rechaza la compra con el mensaje exacto ``Please fill out
        Name and Creditcard.'' y no genera comprobante. Esta es la regla que decide el
        resultado de cada corrida.
  \item \textbf{R3 — El total mostrado = suma de los precios de los productos.} Cuando
        la compra se completa, el monto del comprobante debe ser igual al total del
        carrito.
\end{itemize}

La regla operativa que se desprende de R5 es: la compra tiene éxito \emph{si, y solo
si}, Name y Credit card están ambos rellenados (con el carrito no vacío por R1). Los
campos Country, City, Month y Year no participan en la validación de obligatoriedad. La
Tabla~\ref{tab:derivacion} aplica esta regla a las 8 corridas.

\begingroup
\small
\renewcommand{\arraystretch}{1.2}
\begin{table}[htb]\centering\small
\caption{Derivación del resultado esperado de cada corrida (R1, R3, R5).}\label{tab:derivacion}
\begin{tabularx}{\textwidth}{@{}>{\centering\arraybackslash}p{1.3cm} >{\centering\arraybackslash}p{1.5cm} >{\centering\arraybackslash}p{1.7cm} >{\centering\arraybackslash}p{2.6cm} Y@{}}
\toprule
\textbf{Corrida} & \textbf{Name} & \textbf{Credit card} & \textbf{¿Ambos obligatorios llenos?} & \textbf{Resultado esperado} \\
\midrule
1 & lleno & lleno & Sí & Compra exitosa + comprobante (R3). \\
2 & lleno & vacío & No (falta tarjeta) & Rechazo: ``Please fill out Name and Creditcard.'' \\
3 & lleno & lleno & Sí & Compra exitosa + comprobante (R3). \\
4 & lleno & vacío & No (falta tarjeta) & Rechazo: ``Please fill out Name and Creditcard.'' \\
5 & vacío & lleno & No (falta nombre) & Rechazo: ``Please fill out Name and Creditcard.'' \\
6 & vacío & vacío & No (faltan ambos) & Rechazo: ``Please fill out Name and Creditcard.'' \\
7 & vacío & lleno & No (falta nombre) & Rechazo: ``Please fill out Name and Creditcard.'' \\
8 & vacío & vacío & No (faltan ambos) & Rechazo: ``Please fill out Name and Creditcard.'' \\
\bottomrule
\end{tabularx}
\end{table}
\endgroup

Resultado: \textbf{2 compras exitosas} (corridas 1 y 3) y \textbf{6 rechazos}
(corridas 2, 4, 5, 6, 7 y 8). La corrida \textbf{3} es la más interesante desde el
diseño: Name y Credit card están llenos, pero Country, City y Year quedan vacíos, y aun
así la compra se completa. Esto evidencia que, de los seis campos, el sistema solo
trata como obligatorios dos —el tipo de hallazgo que el arreglo ortogonal expone de
forma económica y que una sola prueba del ``camino feliz'' jamás mostraría—.

\subsubsection*{8.3.7 Casos de prueba formales (IEEE 829)}

Las 8 corridas se documentan como casos de prueba con el formato IEEE 829 de la suite
(ID, Módulo, Prioridad, Técnica de diseño, Precondiciones, Pasos, Datos de prueba,
Resultado esperado y Criterio de referencia). Los identificadores \texttt{CP-AO-01}
\dots{} \texttt{CP-AO-08} no colisionan con ningún ID de la suite. Para mantener la
tabla legible, las tres columnas que son idénticas en los ocho casos se declaran una
sola vez aquí, y la Tabla~\ref{tab:cp-ao} recoge las seis columnas restantes:

\begin{itemize}
  \item \textbf{Módulo:} Orden de compra.
  \item \textbf{Técnica de diseño:} Arreglos ortogonales.
  \item \textbf{Precondiciones:} sesión iniciada y carrito con al menos un producto;
        ventana \emph{Place Order} abierta.
\end{itemize}

\textbf{Criterio de priorización.} Las dos corridas de compra exitosa
(\texttt{CP-AO-01} y \texttt{CP-AO-03}) se marcan \textbf{Alta} porque ejercitan el
camino crítico del negocio —completar una compra—, prioridad máxima según el apartado
3.6 del plan. Las seis corridas de rechazo por R5 se marcan \textbf{Media}, en
coherencia con los casos negativos equivalentes de la suite (\texttt{CP-COM-05} y
\texttt{CP-COM-06}, también Media).

\medskip
\noindent\textbf{Estado de ejecución.} Los ocho casos de la Tabla~\ref{tab:cp-ao} están
\emph{diseñados} en esta sección; su ejecución y resultados se reportan en la
sección 9. No se consignan aquí resultados, capturas ni tiempos de ejecución. El
comportamiento esperado de los campos obligatorios ya está corroborado por los casos
ejecutados \texttt{CP-COM-04}, \texttt{CP-COM-05} y \texttt{CP-COM-06}.

\begingroup
\footnotesize
\setlength{\tabcolsep}{4pt}
\renewcommand{\arraystretch}{1.25}
\begin{longtable}{@{}L{1.5cm} L{1.0cm} L{3.1cm} L{2.8cm} L{3.4cm} L{2.8cm}@{}}
\caption{Casos de prueba de arreglos ortogonales en formato IEEE 829 (\texttt{CP-AO-01}--\texttt{CP-AO-08}). Módulo: Orden de compra; Técnica: Arreglos ortogonales; Precondiciones: sesión iniciada, carrito con al menos un producto y ventana \emph{Place Order} abierta.}\label{tab:cp-ao}\\
\toprule
\textbf{ID} & \textbf{Prior.} & \textbf{Pasos} & \textbf{Datos de prueba} & \textbf{Resultado esperado} & \textbf{Criterio de referencia} \\
\midrule
\endfirsthead
\multicolumn{6}{@{}l}{\footnotesize\emph{Tabla~\ref{tab:cp-ao} (continuación)}}\\
\toprule
\textbf{ID} & \textbf{Prior.} & \textbf{Pasos} & \textbf{Datos de prueba} & \textbf{Resultado esperado} & \textbf{Criterio de referencia} \\
\midrule
\endhead
\midrule
\multicolumn{6}{r@{}}{\footnotesize\emph{continúa en la página siguiente}}\\
\endfoot
\bottomrule
\endlastfoot
CP-AO-01 & Alta & 1.~Rellenar los seis campos con datos válidos. 2.~Pulsar Purchase. & Name: Joshua Mendez; Country: Colombia; City: Bogotá; Credit card: 4111\allowbreak1111\allowbreak1111\allowbreak1111; Month: 12; Year: 2030. & La compra se confirma y se muestra un comprobante con identificador, monto, titular y fecha. El monto es igual al total del carrito (R3). & Convención del dominio / R5: con Name y Credit card completados la compra procede; coherencia con R3 en el comprobante. \\
\midrule
CP-AO-02 & Media & 1.~Rellenar Name, Country y City. 2.~Dejar vacíos Credit card, Month y Year. 3.~Pulsar Purchase. & Name: Joshua Mendez; Country: Colombia; City: Bogotá; Credit card: (vacío); Month: (vacío); Year: (vacío). & El sistema rechaza la compra con el mensaje ``Please fill out Name and Creditcard.'' No genera comprobante ni vacía el carrito. & Convención del dominio / R5: Credit card es obligatorio. \\
\midrule
CP-AO-03 & Alta & 1.~Rellenar Name, Credit card y Month. 2.~Dejar vacíos Country, City y Year. 3.~Pulsar Purchase. & Name: Joshua Mendez; Country: (vacío); City: (vacío); Credit card: 4111\allowbreak1111\allowbreak1111\allowbreak1111; Month: 12; Year: (vacío). & La compra se confirma y se muestra el comprobante (monto = total del carrito, R3), aunque Country, City y Year estén vacíos. Evidencia que esos campos no son obligatorios. & Convención del dominio / R5: solo Name y Credit card son obligatorios; coherencia con R3. \\
\midrule
CP-AO-04 & Media & 1.~Rellenar Name y Year. 2.~Dejar vacíos Country, City, Credit card y Month. 3.~Pulsar Purchase. & Name: Joshua Mendez; Country: (vacío); City: (vacío); Credit card: (vacío); Month: (vacío); Year: 2030. & El sistema rechaza la compra con el mensaje ``Please fill out Name and Creditcard.'' No genera comprobante ni vacía el carrito. & Convención del dominio / R5: Credit card es obligatorio. \\
\midrule
CP-AO-05 & Media & 1.~Rellenar Country, Credit card y Year. 2.~Dejar vacíos Name, City y Month. 3.~Pulsar Purchase. & Name: (vacío); Country: Colombia; City: (vacío); Credit card: 4111\allowbreak1111\allowbreak1111\allowbreak1111; Month: (vacío); Year: 2030. & El sistema rechaza la compra con el mensaje ``Please fill out Name and Creditcard.'' No genera comprobante ni vacía el carrito. & Convención del dominio / R5: Name es obligatorio. \\
\midrule
CP-AO-06 & Media & 1.~Rellenar Country y Month. 2.~Dejar vacíos Name, City, Credit card y Year. 3.~Pulsar Purchase. & Name: (vacío); Country: Colombia; City: (vacío); Credit card: (vacío); Month: 12; Year: (vacío). & El sistema rechaza la compra con el mensaje ``Please fill out Name and Creditcard.'' No genera comprobante ni vacía el carrito. & Convención del dominio / R5: Name y Credit card son obligatorios. \\
\midrule
CP-AO-07 & Media & 1.~Rellenar City y Credit card. 2.~Dejar vacíos Name, Country, Month y Year. 3.~Pulsar Purchase. & Name: (vacío); Country: (vacío); City: Bogotá; Credit card: 4111\allowbreak1111\allowbreak1111\allowbreak1111; Month: (vacío); Year: (vacío). & El sistema rechaza la compra con el mensaje ``Please fill out Name and Creditcard.'' No genera comprobante ni vacía el carrito. & Convención del dominio / R5: Name es obligatorio. \\
\midrule
CP-AO-08 & Media & 1.~Rellenar City, Month y Year. 2.~Dejar vacíos Name, Country y Credit card. 3.~Pulsar Purchase. & Name: (vacío); Country: (vacío); City: Bogotá; Credit card: (vacío); Month: 12; Year: 2030. & El sistema rechaza la compra con el mensaje ``Please fill out Name and Creditcard.'' No genera comprobante ni vacía el carrito. & Convención del dominio / R5: Name y Credit card son obligatorios. \\
\end{longtable}
\endgroup

\subsubsection*{8.3.8 Resumen del bloque de arreglos ortogonales}

\begin{itemize}
  \item \textbf{Técnica:} arreglos ortogonales de Taguchi, la cuarta técnica de caja
        negra que completa la suite.
  \item \textbf{Factores:} los 6 campos de \emph{Place Order}, con 2 niveles cada uno
        (rellenado-válido / vacío).
  \item \textbf{Combinación exhaustiva:} $2^{6}=64$ corridas; \textbf{arreglo elegido:}
        L8 ($2^{7}$), que cubre todas las parejas de niveles (pairwise) en \textbf{8
        corridas}, una reducción del 87,5\,\%.
  \item \textbf{Resultados esperados:} derivados de R1, R3 y R5; 2 compras exitosas
        (\texttt{CP-AO-01}, \texttt{CP-AO-03}) y 6 rechazos con el mensaje ``Please fill
        out Name and Creditcard.''.
  \item \textbf{Hallazgo de diseño:} la corrida 3 muestra que de los 6 campos solo Name
        y Credit card son obligatorios.
  \item \textbf{Estado:} los ocho casos fueron \emph{ejecutados} (11-oct-2026) con
        resultado 8/8 aprobados; el detalle y las evidencias están en la sección 9.3.
\end{itemize}
